\documentclass[12pt]{article}
\usepackage{styles/cwilliams-standard}

% Set up document metadata
\setauthors{Tzu-Chen Chiu, Arjun Sivakumar, Mara Schimmel, Christopher Williams}
\setheaderauthors{Chiu, Sivakumar, Schimmel, Williams}
\settitle{Team Contract - Updated}
\setclass{IPD}

\begin{document}

\maketitlepage

\noindent\textbf{Team Number and Name:} Team 4, Christopher Williams, Mara Schimmel, Arjun Sivakumar, Tzu Chen Chiu

\section{Goals}

\begin{enumerate}
    \item \textbf{What are the goals of the team?}

    As a team we will create a transformer-based neural network to decode quantum codes. If time permits, we will compare our system to one developed by Google. Our main goal is to see how it performs on the Surface Code and other commonly used quantum codes.

    \item \textbf{What are the personal goals of each individual team member?}
    \begin{enumerate}
        \item Christopher Williams: Learn from industry professionals and become familiar with Quantum Computing.
        \item Arjun Sivakumar: My goal out of this project is to mainly gain some experience in industry-standard projects and programming, while also learning about quantum computing and error correction.
        \item Mara Schimmel: Learn from industry professionals and get more experience working outside of personal projects.
        \item Tzu Chen: Gain industry experience. Learn this new area directly with experts from industry.
    \end{enumerate}

    \item \textbf{What happens if all of you decide you want to get an A grade, but because of time constraints, one person decides that a B will be acceptable?}

    As a team, if one or two people are having difficulties and are time constrained, the rest of the team will provide help and assistance. If this is not enough, we will reevaluate our distribution of responsibilities.

    \item \textbf{What happens if one or more team members want to continue the project, release the project as open source or give the project to the company mentor (if any) after the course finishes?}

    After the project concludes, as a group we are content with discussing what our intentions are for the project at a later date or near the end of the project timeline.
\end{enumerate}

\section{Communication}

\begin{enumerate}
    \setcounter{enumi}{4}
    \item \textbf{How will we communicate (remote video call, text, email, chat (slack, discord, etc))?}
    \begin{enumerate}
        \item Discord and group chats are the chosen method of communication.
    \end{enumerate}

    \item \textbf{What are acceptable times to contact other team members by text? By chat? Etc?}
    \begin{enumerate}
        \item Any time works for our team by text or chat.
    \end{enumerate}

    \item \textbf{How quickly must responses to team member questions/contacts be answered?}
    \begin{enumerate}
        \item Depending on the level of priority, response time will vary. If a team member is providing an update, an acknowledgement of any kind is appreciated.
    \end{enumerate}

    \item \textbf{Where will we meet to discuss the project (classroom, library, remote video)? It's ok to meet online, but when you do, everyone must be focused on the meeting not distracted with some other task or game you are on.}
    \begin{enumerate}
        \item If meeting to discuss large project goals and decisions, team will meet in person at AB1. Otherwise, meetings can be scheduled online or in person.
    \end{enumerate}

    \item \textbf{How often will we meet (\underline{minimum once a week})? This is common for successful teams.}
    \begin{enumerate}
        \item As a team we will meet minimum once a week, and other times will be decided upon when necessary.
    \end{enumerate}

    \item \textbf{When will we meet (\underline{the exact day and time})? These should be specific and sacrosanct. This is a common trait of successful teams. Blowing off meetings will be a big issue with the instructor. Not scheduling enough meetings as needed during times of distress will also be an issue.}
    \begin{enumerate}
        \item As a team, we will definitely meet on Wednesdays at 1:30 pm in AB1. Other meeting times will be decided upon, and if necessary we will meet at 1:30pm on Mondays.
    \end{enumerate}
\end{enumerate}

\section{Work Expectations}

\begin{enumerate}
    \setcounter{enumi}{10}
    \item \textbf{How many hours will each team member be expected to work? Keep in mind that the instructional staff require each member to work \underline{a minimum of 10 hours a week EVERY week} on the project. Once the Sprints begin these hours do not include meetings or generating project documentation or deliverables. Be reasonable, but if the work expectations don't match the aggressiveness of the MVP, there may be problems with your teammates and the instructor for the course.}
    \begin{enumerate}
        \item During the learning process (the first month of the project), each team member will dedicate 10-15 hours a week to the project. Once the learning phase is over, each team member will dedicate 15-20 hours a week towards the project. If their responsibilities are completed, they do not need to meet the time goal.
    \end{enumerate}

    \item \textbf{How will you handle it if the project is behind schedule? Do you expect team members to work all hours towards a deadline if they are behind? Do you expect others to help?}
    \begin{enumerate}
        \item If one or two team members are behind, the other team members will assist in their responsibilities, assisting them to complete what is asked of them. If the entire team is falling behind, we will reevaluate the responsibilities and goals of the team.
    \end{enumerate}

    \item \textbf{What are the expected work hours for each team member? Do some team members work and need to complete all of their hours on the weekend? Do some team members need to work a few hours each evening or morning? Keep in mind that the team members' schedules may affect other team members if they are waiting for a work product.}
    \begin{enumerate}
        \item Team members will complete their work at their preferred times. Team members will complete their hours by the end of the week and when they see fit.
    \end{enumerate}

    \item \textbf{How will you use in-class time, if given any? Remember attendance is expected.}
    \begin{enumerate}
        \item Team members will work on their individual responsibilities and discuss topics that the group is working on.
    \end{enumerate}

    \item \textbf{How will you handle semester breaks like Fall Break...what do you expect for this time? A start-up would die if everyone took time off. What will your team do?}
    \begin{enumerate}
        \item The team will work ahead and complete responsibilities before fall break so that all team members can enjoy their deserved times off.
    \end{enumerate}

    \item \textbf{Is it acceptable for one or two team members to choose to do extra work than the others in order to get the whole team an A?}
    \begin{enumerate}
        \item If team members ``pick up the slack'' of others at an irregular interval, that is acceptable. If it is a consistent situation that prevents a team member from completing their responsibilities, the team will redistribute said responsibilities.
    \end{enumerate}

    \item \textbf{What will you do if one or more team members are not doing their share of the work?}
    \begin{enumerate}
        \item Discuss with the team member(s) about their difficulties, and decide if the team as whole can assist with their responsibilities.
    \end{enumerate}

    \item \textbf{What will you do if one or more team members are individually deciding to add tasks, change their assigned tasks or perform additional work?}
    \begin{enumerate}
        \item Discuss with the team about the change or additions and confirm it aligns with the goals of the team.
    \end{enumerate}

    \item \textbf{Is it acceptable for one team member to do work assigned to another team member and if so under what circumstances?}
    \begin{enumerate}
        \item If a team member is struggling, other team members may ask about assisting with their responsibilities and helping out. Discuss before doing work of other team members.
    \end{enumerate}

    \item \textbf{What will happen if someone does not follow through on a commitment (e.g., missing a deadline, not showing up to meetings)?}
    \begin{enumerate}
        \item The team will check in with that member and assure the team member is doing well. If the team member is doing well, a discussion about responsibilities will be had as a team if the issue is consistent, if the issue cannot be resolved within the team we will seek assistance from the professor.
    \end{enumerate}

    \item \textbf{How will you deal with different work habits of individual team members (e.g., some people like to get assignments done as early as possible; others like to work under the pressure of a deadline)?}
    \begin{enumerate}
        \item The team will be flexible and allow team members to finish their work at their own pace, but still collaborate and discuss issues as they arise and get help from other team members. If a team member's work relies on another team member's progress, work times should be flexible and conscious of how their responsibilities affect other team members.
    \end{enumerate}

    \item \textbf{How will you handle situations that require team members to work together when they have different work hours?}

    Since we are all full time students, we are free at night. Or maybe we could work before/after class.

    \item \textbf{What are the individual team member responsibilities if another team member is over or under reporting their portion of any tasks they worked on? If they are over or under reporting the time required to complete tasks? If there are any other violations of the Honor code?}
    \begin{enumerate}
        \item The team will discuss with the team members that are under or over reporting about their progress and work towards a solution that is more accurate of the effort. If necessary, we will discuss it with the professor.
    \end{enumerate}
\end{enumerate}

\section{Team Member Roles}

\begin{enumerate}
    \setcounter{enumi}{23}
    \item \textbf{What will the roles of the team members be? (Team lead, notetaker, manager, member etc.) At a minimum you must have a Team Lead. The team lead will rotate every 2 weeks. Any other positions you assign may change with the Team lead or stay the same. Please list who will be the team lead for each section. Everyone must take a turn.}
    \begin{enumerate}
        \item Sprint 1
        \begin{enumerate}
            \item[i.] Team lead: Arjun
            \item[ii.] Notetaker: Chris
        \end{enumerate}
        \item Sprint 2
        \begin{enumerate}
            \item[i.] Team lead: Tzu Chen
            \item[ii.] Notetaker: Chris
        \end{enumerate}
        \item Sprint 3
        \begin{enumerate}
            \item[i.] Team lead: Chris
            \item[ii.] Notetaker: Mara
        \end{enumerate}
        \item Sprint 4
        \begin{enumerate}
            \item[i.] Team lead: Mara
            \item[ii.] Notetaker: Chris
        \end{enumerate}
        \item Sprint 5
        \begin{enumerate}
            \item[i.] Team lead: Chris
            \item[ii.] Notetaker: Arjun
        \end{enumerate}
        \item Until end of semester (you may leave blank if you have less than 6 team members)
    \end{enumerate}

    \item \textbf{Who in our team will submit team assignments when they are due?}
    \begin{enumerate}
        \item The team lead will submit team assignments during their sprint.
    \end{enumerate}

    \item \textbf{Who in our team will reach out to the instructional staff when there is a question or concern about the project? This should probably be the Team Lead.}
    \begin{enumerate}
        \item The team lead will reach out to the instructional staff.
    \end{enumerate}

    \item \textbf{Where will meeting minutes be kept and by whom? This should probably be the notetaker if you have one.}
    \begin{enumerate}
        \item Meeting minutes will be kept in a shared google doc, added to by the notetaker for each meeting.
    \end{enumerate}
\end{enumerate}

\section{Work Products/Code Review}

\begin{enumerate}
    \setcounter{enumi}{27}
    \item \textbf{How will the work be reviewed? You should definitely have a code review process.}
    \begin{enumerate}
        \item Choose a specific day to perform code review per week, where the available member picks up the other member's code and does a code review for it, along with relevant constructive criticisms.
    \end{enumerate}

    \item \textbf{What happens if people have different opinions on the quality or acceptability of the work?}
    \begin{enumerate}
        \item We can have a chat about it over text, call or in person meet up, presenting our perspectives and then try to reach a consensus. If the issue or disagreement is larger and affects the entire team, then we can have a chat as a group or even take it up with the relevant professor, who could provide more insight on correct code quality.
    \end{enumerate}

    \item \textbf{Who ``owns'' any work produced (in terms of code ownership not IP)? Who can make changes to other team members' work and under what circumstances?}
    \begin{enumerate}
        \item Each member has complete ownership of their own code they have written, and if large-scale changes are to be made to the other member's code, the member must be made aware of this decision, and they have the right to deny.
    \end{enumerate}
\end{enumerate}

\section{Decisions/Conflict/Difficult Conversations}

\begin{enumerate}
    \setcounter{enumi}{30}
    \item \textbf{How will decisions be made? (For example: Majority vote? Does the Team Lead have the final say?)}

    We will discuss as a team and try to agree. If we cannot, we will use a majority vote.

    \item \textbf{How will you handle a team member who is being difficult or not being respectful (and what does difficult and respectful mean) or not letting all voices be heard?}

    We will first remind the person directly and politely, and if it continues, we will ask the professor for help.

    \item \textbf{How will individual teammates deal with difficulty with another individual teammate?}

    They could have a private meeting at first. If the problem continues, they will bring it to the team.

    \item \textbf{How will you handle unforeseen conflicts that arise?}

    Discuss with the group immediately.

    \item \textbf{Any other issues you feel need to be addressed.}
    \begin{enumerate}
        \item None!
    \end{enumerate}
\end{enumerate}

\section*{Signatures}

\begin{figure}[H]
    \centering
    \includegraphics[width=0.75\textwidth]{images/team_contract_signautres.png}
\end{figure}

\section*{List of Changes}

\begin{itemize}
    \item Removed Dai Chen from contract (still included in signatures section).
    \item Adjusted Sprint assignments.
    \item Small update to Goals section.
\end{itemize}

\end{document}
